% =====================================================================
%  Chapter 6 — The Method, Step by Step (new Part III)
% =====================================================================
\part{The Method, Step by Step}
\label{part:method}

\chapter{The Shape of the Argument}
\label{ch:shape}

The previous part set out the problem and the walls. This part does the thing
the owner of the claim would most want done and the coverage most conspicuously
failed to do: it sets out \emph{the method itself}, step by step, in the terms a
mathematician can follow, so that a reader can arrive at the result rather than
be told about it.

There is an important qualification, and it belongs at the front. The account
below is assembled from two sources, and they are not of equal quality.

\begin{quote}
\textbf{The remark} gives a twenty-item timeline that its own author calls ``a
rough timeline Claude provided'' and, of the parts he did not personally
develop, ``sort of paraphrasing a `slop' summary, which might make zero sense''.
Several of those entries are one sentence long and are addressed to no reader in
particular.

\textbf{The Lean development} is the authoritative version. It contains, in
Lean, the actual statements of the step-by-step mathematics: the chain
operators, the integrality conditions on the chain, the descent-tower
conditions, the seesaw, the Behrend set, the rank budget. Those statements are
precise, and they are what the arguments of this part are anchored to.
\end{quote}

So the method is presented here in a specific way: the mathematical content is
taken from the formal statements wherever they exist, because those are
precise; the narrative connective tissue is reconstructed from the remark, and
is flagged as reconstruction wherever the remark does not state it outright.
Chapter~\ref{ch:map} gives the complete step-to-declaration map, so a reader
who wants to check any sentence of this part can do so in one move.

\section{The claim in one sentence}\begin{quote}
The Riemann hypothesis is the statement that a certain real-valued function
built from the ordinates of the zeros of $\zeta$ --- the Weil expression $W$ ---
is non-negative for every real parameter. This argument manufactures, by
iterating theta lifts, the arithmetic multiplicity that the classical
function-field proof gets for free; it converts that multiplicity into a
positivity statement for $W$ by a determinant; it converts the positivity
statement into a bound on the zeros by the Siegel--Weil formula and a
test-function argument; and it kills the leftover cohomology by twisting,
which is where the open part of the problem sits.
\end{quote}

That is the whole shape. What follows is the nine steps that make it up, and
what each one costs.

\section{Nine steps}\begin{center}
\small
\setlength{\extrarowheight}{-1pt}
\begin{tabular}{@{}rl>{\RaggedRight\arraybackslash\hyphenpenalty=10000\exhyphenpenalty=10000}p{0.56\textwidth}@{}}
\toprule
\textbf{Step} & \textbf{Remark item} & \textbf{What it does} \\
\midrule
1 & 1, 4 & Establish the ancestor: what Stepanov's method does with
multiplicity, and what has to replace it over the integers. \\
2 & 2, 9, 10 & Build the divisibility engine: a CM curve, an invariant
derivation, and a prime-power divisibility that pays for level and
cancellation. \\
3 & 3, 5 & Produce the identity $qN = Mk + R$ and the sign-only Weil
argument: a determinant that is a Gram matrix, and a remainder with
square-root quality. \\
4 & 8 & Chain the theta lifts: $U(X_0) \to \cdots \to U(X_m)$, each step
winning one unit of multiplicity, at a price. \\
5 & 14 & Compute the price: the budget that each lift spends against the
multiplicity it buys. This is where the two-thirds ceiling appears. \\
6 & 12 & Apply the Rallis seesaw, so that the period of the iterated lift
becomes a Bessel period, independent of the Hecke characters. \\
7 & 11, 19 & Average over the Hecke characters: Behrend's pigeonhole
construction, the variance bounds, and the equidistribution of translates
that produces the extra averaging. \\
8 & 16, 18 & Secure the automorphy and the arithmetic growth: the relative
trace formula, the sixfold cover of $\mathrm{GL}_5$, the Poincar\'e series,
and the Iwasawa tower. \\
9 & 17 & Kill the surviving cohomology: twist by characters, and confront
what survives. \\
\bottomrule
\end{tabular}
\end{center}

Chapters~\ref{ch:step1} to~\ref{ch:vanishing} take these in order.

\section{What ``multiplicity'' means here, and why it is the crux}Before the steps, one idea has to be stated precisely, because everything in
the argument is about it.

In the function-field proof of the Riemann hypothesis for curves over finite
fields, the argument produces a divisor $D$ on the curve with a point of high
multiplicity, and the dimension of the space of functions with poles bounded by
$D$ is large enough to force a contradiction with the count of points. The
count of points over $\mathbb{F}_q$ is $q + O(\sqrt q)$; the dimension of the
function space is a geometric quantity; and the *product* of the two is what
creates the bound. Formally the ambient object is a product of a geometric
factor and an arithmetic factor, and multiplicity is what the two factors
multiply.

Over the integers the arithmetic factor is $\operatorname{Spec}\mathbb{Z}$ and
the product is one-dimensional. The dimension count has nothing to multiply
against.

\begin{quote}
\textbf{The idea of the whole method.} Find an operation that behaves, for the
purposes of the dimension count, like the missing arithmetic factor: an
operation that multiplies the size of the relevant space by a controllable
factor. The theta lift is that operation. Each application of it is claimed to
multiply the space by one further unit of multiplicity, and the chain
$U(X_0) \to \cdots \to U(X_m)$ is the attempt to apply it enough times.
\end{quote}

That is the sentence the remark offers as its one-line summary, and it is
genuinely the crux:

\begin{quote}
\small
``for the von Mangoldt sum of length $X$, Claude's suggested $\sqrt{\ }$ length of
chain of theta lifts, in view of each mimicking multiplicity, was $X$. (And yet,
what else could it be?)''
\end{quote}

The parenthetical is the author recording that he cannot explain why the
natural length should be $\sqrt{X}$. It is worth being careful: $\sqrt{X}$ is
the right answer for a reason that can be stated. If each lift multiplies the
space by a factor $\lambda$, then after $m$ lifts the space has grown by
$\lambda^m$, and to dominate a sum of length $X$ one needs $\lambda^m \ge X$;
with $\lambda$ a fixed constant, $m \ge \log X / \log \lambda$. A chain of
length $\sqrt{X}$ is what one gets if $\lambda$ is itself allowed to grow with
$X$ --- that is, if each lift buys a *multiplicative* rather than an
*additive* amount of multiplicity. That is a much stronger claim than ``each
lift wins multiplicity $m$'' and it is the claim that carries the whole
argument.

So the reader should hold two quantities apart:

\begin{description}[style=nextline,leftmargin=0pt,labelwidth=0pt,labelsep=0pt,font=\sffamily\bfseries\small,
  itemsep=6pt]
\item[\textbf{Multiplicity gained}] How much the function space grows per lift.
This is what the lift chain is supposed to provide.
\item[\textbf{Cost}] What the lift spends --- denominators, levels, the size of
the parameters --- and therefore how much of the analytic budget is gone
before the argument reaches the point where the bound is read off.
\end{description}

Chapter~\ref{ch:budget} computes the second, and the result is the
two-thirds figure.

\section{How the argument is anchored in the source}Every step above has a corresponding statement in the submitted development.
The mapping is:

\begin{center}
\begin{tabular}{@{}>{\RaggedRight\arraybackslash\hyphenpenalty=10000\exhyphenpenalty=10000}p{0.22\textwidth}>{\RaggedRight\arraybackslash\hyphenpenalty=10000\exhyphenpenalty=10000}p{0.66\textwidth}@{}}
\toprule
\textbf{Step} & \textbf{Where the mathematics lives in the development} \\
\midrule
1 & \code{Development/WeilBound/RepresentationTheoryGL/SiegelWeilProof/} and
\code{.../ThetaLifts/LeviDeterminants/} \\
2 & \code{Formalism/SharpRestriction/Divisibility/} \\
3 & \code{Development/WeilBound/RepresentationTheoryGL/RegularisedSiegelWeil/} \\
4 & \code{Development/WeilBound/ThetaLifts/} and
\code{Formalism/ThetaLiftIndependence/} \\
5 & \code{Formalism/SharpRestriction/Divisibility/ChainOperators/} \\
6 & \code{Formalism/SharpRestriction/Divisibility/FunctionModelRallisFromInputs.lean} \\
7 & \code{Formalism/BehrendSets/} \\
8 & \code{Development/WeilBound/CoreStatements/AutomorphicKernels/} \\
9 & \code{Development/AbelianVarietyTheory/Definitions/Foundations/FinalTheorems.lean} \\
\bottomrule
\end{tabular}
\end{center}

Chapter~\ref{ch:map} gives this in full, with the specific declarations.

\section{What this part does and does not claim}It claims to explain the method as the author describes it and as the
development formalises it, in enough detail that a reader with a mathematics
background can follow the argument's structure and see exactly where each
step's weight comes from.

It does not claim that the argument is correct. Several of the steps are open
problems in the literature as of this writing, one step is known to the author
himself to have a ceiling, and the formalisation has never been audited by an
outside party. What this part does is make the claim \emph{legible}, which is
the precondition for anyone judging it --- including the reader.

\begin{keytakeaways}
\item The method is one sentence: manufacture the missing arithmetic
multiplicity by iterating theta lifts, convert it to a Weil positivity by a
determinant, convert that to a bound on the zeros, and kill the leftover
cohomology by twisting.
\item It has nine steps, and each step has a precise statement in the
submitted development.
\item Multiplicity gained and cost paid are two different quantities; the
argument's difficulty is concentrated in their ratio.
\item The chain length $\sqrt{X}$ is what one gets if each lift buys a
multiplicative rather than an additive amount of multiplicity --- a strong
claim, and the one that carries the argument.
\end{keytakeaways}
